\chapter{Główne twierdzenia i algorytmy}

Ten rozdział został opracowany na podstawie następujących prac i monografii \cite{eshwar2025,borkar2008,puterman2014,Kallenberg2022}.

\section*{Twierdzenie o redukcji polityk}
Pierwszym krokiem w analizie jest pokazanie, że dowolną politykę można zredukować do prostszej formy: decyzja w kroku $t=0$ może być dowolna, natomiast od kroku $t\ge 1$ wystarczy rozważać politykę stacjonarną.

\begin{theorem}[Redukcja polityk]
\label{thm:policy-reduction}
Dla MDP o skończonych zbiorach stanów i akcji, dla dowolnej polityki  $\phi = (\phi_0, \phi_1, \ldots)$ istnieje reguła pierwszego kroku $\mu$ oraz stacjonarna polityka $\phi_s$ takie, że
\[
  J_{\phi}^{\alpha,\beta}(s) = J_{(\mu,\phi_s)}^{\alpha,\beta}(s), \qquad \forall s \in S.
\]

\end{theorem}

Z Twierdzenia~\ref{thm:policy-reduction} wynika, że znalezienie optymalnej polityki sprowadza się do analizy pary składającej się z reguły decyzyjnej pierwszego kroku $\mu$ oraz stacjonarnej polityki $\phi_s$: $(\mu,\phi_s)$.

\begin{proof}
Dla dowolnej polityki $\phi = (\phi_0, \phi_1, \ldots)$ mamy
\[
  J_{\phi}^{\alpha,\beta}(s) = \mathbb{E}\!\left[ \sum_{t=0}^{\infty} d(t)\, r(s_t, a_t) \;\middle|\; s_0 = s \right],
\]
gdzie $d(0) = 1$ oraz $d(t) = \alpha\,\beta^{t}$ dla $t \ge 1$.

Rozwijając wartość oczekiwaną otrzymujemy
\[
  J_{\phi}^{\alpha,\beta}(s) = \mathbb{E}\!\left[ r(s_0,a_0) + \sum_{t=1}^{\infty} \alpha \beta^{t}\, r(s_t,a_t) \;\middle|\; s_0 = s \right].
\]

Warunkując po pierwszej akcji $a_0 \sim \phi_0(\cdot \mid s)$ dostajemy
\[
  J_{\phi}^{\alpha,\beta}(s) = \mathbb{E}_{a_0 \sim \phi_0(\cdot \mid s)}\!\left[ r(s,a_0) + \,\mathbb{E}\!\left[ \sum_{t=1}^{\infty} \alpha \beta^{t}\, r(s_t,a_t) \;\middle|\; s_0 = s, a_0 \right] \right].
\]

Kolejne warunkowanie względem $s' \sim q(\cdot \mid s,a_0)$ daje
\[
  J_{\phi}^{\alpha,\beta}(s) = \mathbb{E}_{\substack{a_0 \sim \phi_0(\cdot \mid s) \\ s' \sim q(\cdot \mid s,a_0)}}\!\left[ r(s,a_0) + \alpha\beta\,\mathbb{E}\!\left[ \sum_{t=1}^{\infty} \beta^{t-1}\, r(s_t,a_t) \;\middle|\; s_1=s' \right] \right].
\]

W dalszym ciągu (dla przejrzystości zapisu) warunkujemy jedynie po $s_1=s'$; formalnie można równoważnie warunkować po całej historii $(s_0=s, a_0, s_1=s')$.

Pozostała suma dyskontowana ma postać wykładniczą, więc
\[
  \mathbb{E}\!\left[ \sum_{t=1}^{\infty} \beta^{t-1}\, r(s_t,a_t) \;\middle|\; s_1 = s' \right] = \,J_{\overrightarrow{\phi}}^{\beta}(s_1),
\]
gdzie $\overrightarrow{\phi} = (\phi_1, \phi_2, \ldots)$ i $ J^{\beta}_{\overrightarrow{\phi}}$ odpowiada klasycznemu dyskontowaniu wykładniczemu. Z Twierdzenia 9.12 w \cite{Kallenberg2022} wynika, że dla ustalonego stanu początkowego $s\in S$ i dla każdej polityki $\overrightarrow{\phi}$ istnieje polityka stacjonarna $\phi_s$ taka, że
$$\sum_{s'} J^\beta_{\overrightarrow{\phi}}(s')q(s'|s,\phi_0(s))= 
\sum_{s'}   J^\beta_{\phi_s} (s')q(s'|s,\phi_0(s)), $$
gdzie $q(s'|s,\phi_0(s))=\sum\limits_a  q(s'|s,a) \phi_0(a|s).$

Podstawiając do poprzedniego równania, otrzymujemy
\[
  J_{\phi}^{\alpha,\beta}(s) = \sum_{a}\phi_{0}(a \mid s)\!\left[ r(s,a) + \alpha\,\beta\, \sum_{s'} q(s' \mid s,a) J^{\beta}_{\phi_s}(s') \right].
\]

Wybierając $\mu = \phi_0$ jako regułę pierwszego kroku, uzyskujemy żądaną równość
\[
  J_{\phi}^{\alpha,\beta}(s) = J_{(\mu,\phi_s)}^{\alpha,\beta}(s), 
\]
co kończy dowód.
\end{proof}

\section{Ocena polityki}
W tym rozdziale przedstawiamy algorytm wyliczenia jednokrokowej niestacjonarnej polityki $(\mu, \phi_s)$ w środowisku MDP z dyskontowaniem quasi-hiperbolicznym. W~przeciwieństwie do dyskontowania wykładniczego, funkcja wartości $J_{\mu, \phi_s}^{\alpha, \beta}$ nie posiada prostej formy rekurencyjnej.
\subsection{Wyprowadzenie wzoru do oceny polityki}

Z Twierdzenia~\ref{thm:policy-reduction}:
\begin{equation}
J^{\alpha, \beta}_{\mu, \phi_s}(s) = \sum_a \mu(a\mid s) \left[ r(s, a) + \alpha\beta \sum_{s'} q(s'\mid s, a) \, J^{\beta}_{\phi_s}(s') \right].
\label{eq:policy-eval-start}
\end{equation}

Celem jest wyrażenie   $J^{\alpha, \beta}_{\mu, \phi_s}(s)$ przez $J^{\alpha, \beta}_{\phi_s}(s')$.
\par\noindent\textbf{Definicje dla polityki stacjonarnej.}\par

Dyskontowanie wykładnicze:
\begin{equation}
J^{\beta}_{\phi_s}(s') = \mathbb{E}_{\phi_s} [r(s', \cdot)] + \beta \mathbb{E}_{\phi_s, q} [J^{\beta}_{\phi_s}(s'') \mid s'],
\label{eq:exp-value}
\end{equation}
 gdzie $\mathbb{E}_{\phi_s} [r(s', \cdot)] = \sum_{a'} \phi_s(a'\mid s') \, r(s', a')$ \\ $\mathbb{E}_{\phi_s,q}[J^{\beta}_{\phi_s}(s'')\mid s'] = \sum_{a'} \phi_s(a'\mid s') \sum_{s''} q(s''\mid s', a') J^{\beta}_{\phi_s}(s'')$.

Dyskontowanie quasi-hiperboliczne:
\begin{equation}
J^{\alpha, \beta}_{\phi_s}(s') = \mathbb{E}_{\phi_s} [r(s', \cdot)] + \alpha \beta \mathbb{E}_{\phi_s, q} [J^{\beta}_{\phi_s}(s'') \mid s'].
\label{eq:qh-value-stat}
\end{equation}

\par\noindent\textbf{Obliczenia:}\par

Pomnóżmy równanie~\eqref{eq:exp-value} przez $\alpha \beta$:
\begin{equation}
\alpha \beta J^{\beta}_{\phi_s}(s') = \alpha \beta \mathbb{E}_{\phi_s} [r(s', \cdot)] + \alpha \beta^2 \mathbb{E}_{\phi_s, q} [J^{\beta}_{\phi_s}(s'') \mid s'].
\label{eq:mult-ab}
\end{equation}

Pomnóżmy równanie~\eqref{eq:qh-value-stat} przez $\beta$:
\begin{equation}
\beta J^{\alpha, \beta}_{\phi_s}(s') = \beta \mathbb{E}_{\phi_s} [r(s', \cdot)] + \alpha \beta^2 \mathbb{E}_{\phi_s, q} [J^{\beta}_{\phi_s}(s'') \mid s'].
\label{eq:mult-b}
\end{equation}

Odejmując równanie~\eqref{eq:mult-b} od równania~\eqref{eq:mult-ab} otrzymujemy:
\begin{equation}
\alpha \beta J^{\beta}_{\phi_s}(s') - \beta J^{\alpha, \beta}_{\phi_s}(s') = \alpha \beta \mathbb{E}_{\phi_s} [r(s', \cdot)] - \beta \mathbb{E}_{\phi_s} [r(s', \cdot)].
\end{equation}

oraz:
\begin{equation}
\alpha \beta J^{\beta}_{\phi_s}(s') = \beta J^{\alpha, \beta}_{\phi_s}(s') - (1 - \alpha) \beta \mathbb{E}_{\phi_s} [r(s', \cdot)].
\label{eq:isolated}
\end{equation}

Podstawiając równanie~\eqref{eq:isolated} do równania~\eqref{eq:policy-eval-start} dostajemy:
\begin{align}
J^{\alpha, \beta}_{\mu, \phi_s}(s) &= \sum_a \mu(a\mid s) \left[ r(s, a) + \sum_{s'} q(s'\mid s, a) \left[ -(1 - \alpha)\beta \sum_{a'} \phi_s(a'\mid s') r(s', a') + \beta J^{\alpha, \beta}_{\phi_s}(s') \right] \right].
\label{eq:final-policy-eval}
\end{align}
W szczególności, jeśli $\mu$ zastąpimy przez $\phi_s$ otrzymujemy:
\begin{align}
J^{\alpha, \beta}_{\phi_s}(s) &= \sum_a \phi_s(a\mid s) \left[ r(s, a) + \sum_{s'} q(s'\mid s, a) \left[ -(1 - \alpha)\beta \sum_{a'} \phi_s(a'\mid s') r(s', a') + \beta J^{\alpha, \beta}_{\phi_s}(s') \right] \right].
\label{eq:final-policy-eval-stationary}
\end{align}

\subsection{Algorytm 1: Obliczenie funkcji wartości dla dyskontowania QH}

\begin{algorithm}[H] \
\caption{Obliczenie funkcji wartości dla dyskontowania QH \cite{eshwar2025}}
\label{alg:policy-eval}
\begin{algorithmic}[1]
\STATE \textbf{Wymagania:} Kroki uczenia $(\eta_n)$, $(\theta_n)$; współczynniki dyskontowania $\alpha$, $\beta$; polityka próbkowania $\nu$; polityki oceniane $\mu$, $\phi_s$
\STATE \textbf{Inicjalizacja:} $J_0, W_0 \in \mathbb{R}^{|S|}$
\FOR{$n = 0, 1, 2, \dots$}
\FOR{$s \in S$}
\STATE Próbkuj $a \sim \nu(\cdot|s)$
\STATE Zaobserwuj nagrodę $r(s, a)$ oraz stan kolejny $s' \sim q(\cdot|s, a)$
\STATE Próbkuj $a' \sim \phi_s(\cdot|s')$
\STATE Zaobserwuj nagrodę $r(s', a')$
\STATE Oblicz cel TD: $r^{\text{target}}_n(s) \leftarrow r(s, a) - (1 - \alpha)\beta \, r(s', a') + \beta W_n(s')$
\STATE Zaktualizuj $W$: $W_{n+1}(s) \leftarrow W_n(s) + \eta_n \left[ \frac{\phi_s(a|s)}{\nu(a|s)} r^{\text{target}}_n(s) - W_n(s) \right]$
\STATE Zaktualizuj $J$: $J_{n+1}(s) \leftarrow J_n(s) + \theta_n \left[ \frac{\mu(a|s)}{\nu(a|s)} r^{\text{target}}_n(s) - J_n(s) \right]$
\ENDFOR
\ENDFOR
\STATE \textbf{Zwraca:} $J_n$ zbieżne do $J^{\alpha, \beta}_{\mu, \phi_s}$; $W_n$ zbieżne do $J^{\alpha, \beta}_{\phi_s}$
\end{algorithmic}
\end{algorithm}


Pierwszy algorytm pozwala na obliczenie funkcji wartości dla danej polityki bez znajomości modelu środowiska. Próbki danych generowane są przez politykę próbkowania $\nu$, a korekta off-policy jest realizowana przez wagi $\frac{\mu(a\mid s)}{\nu(a\mid s)}$ oraz $\frac{\phi_s(a\mid s)}{\nu(a\mid s)}$. W szczególności wymagane jest pokrycie: $\nu(a\mid s) > 0$ dla każdej akcji, która może zostać wybrana przez $\mu$ lub $\phi_s$ w danym stanie. W tej pracy będziemy stosować $\nu$ o rozkładzie jednostajnym.
\subsection{Mechanizm aktualizacji wartości i cel TD}

Kluczowym elementem Algorytmu 1 jest specyficzna konstrukcja celu TD (ang. \textit{Temporal Difference target}) oraz wykorzystanie dwóch skal czasowych do jednoczesnej aproksymacji pomocniczej funkcji wartości $W$ (dla dyskontowania wykładniczego) i docelowej funkcji wartości $J$ (dla dyskontowania quasi-hiperbolicznego).

\subsubsection*{1. Cel TD}

W standardowym algorytmie Q-learning cel TD opiera się na wyrażeniu $r + \beta \max Q$. W przypadku dyskontowania quasi-hiperbolicznego, ze względu na strukturę równania rekurencyjnego (3.8), aktualzacja celu musi uwzględniać korektę nagrody w następnym kroku. Cel $r_n^{\text{target}}$ w $n$-tej iteracji definiujemy jako:

\begin{equation}
    r_n^{\text{target}}(s) = r(s,a) - (1-\alpha)\beta r(s',a') + \beta W_n(s')
    \label{eq:td_target}
\end{equation}

\noindent gdzie:
\begin{itemize}
    \item $r(s,a)$ to natychmiastowa nagroda w obecnym stanie,
    \item $-(1-\alpha)\beta r(s',a')$ to człon korygujący wynikający z definicji operatora Bellmana dla QH Lematu \ref{thm:bellman-contraction}, uwzględniający różnicę w postrzeganiu nagrody przez ,,ja'' teraźniejsze i przyszłe,
    \item $W_n(s')$ to przybliżona funkcja wartości wykładniczej w następnym stanie.
\end{itemize}

\subsubsection*{2. Korekta Off-Policy (Importance Sampling)}

Skoro próbkowanie akcji odbywa się według polityki behawioralnej $v$ (Rozkład jednostajny), to chcemy ocenić politykę złożoną z $\mu$ oraz $\phi_s$, stosujemy wagi wiarygodności (ang. \textit{importance sampling ratios}):

\begin{itemize}
    \item Do aktualizacji $W$ (związanej z polityką $\phi_s$):
    \begin{equation}
        \rho_n^W = \frac{\phi_s(a|s)}{v(a|s)}
    \end{equation}
    \item Dla aktualizacji $J$ (związanej z polityką $\mu$):
    \begin{equation}
        \rho_n^J = \frac{\mu(a|s)}{v(a|s)}
    \end{equation}
\end{itemize}

\subsubsection*{3. Równania aktualizacji (Dwie skale czasowe)}

Aktualizacja parametrów odbywa się w dwóch skalach czasowych, co gwarantuje stabilność zbieżności:

\paragraph{Szybka skala (Aproksymacja $W$):} Funkcja $W_n$ jest aktualizowana z krokiem $\eta_n$ i dąży ona do wartości $J_{\phi_s}^{\alpha, \beta}$.

\begin{equation}
    W_{n+1}(s) \leftarrow W_{n}(s) + \eta_{n} \left[ \rho_n^W r_{n}^{\text{target}}(s) - W_{n}(s) \right].
    \label{eq:update_W}
\end{equation}

\paragraph{Wolna skala (Aproksymacja $J$):} Funkcja $J_n$, będąca naszym głównym celem (wartość QH dla pary $(\mu, \phi_s)$), jest aktualizowana z mniejszym krokiem $\theta_n$. Wykorzystuje ona ten sam cel $r_n^{\text{target}}$, ale skorygowany wagą dla polityki pierwszego kroku $\mu$.

\begin{equation}
    J_{n+1}(s) \leftarrow J_{n}(s) + \theta_{n} \\ \left[ \rho_n^J r_{n}^{\text{target}}(s) - J_{n}(s) \right].
    \label{eq:update_J}
\end{equation}

Warunek $\frac{\theta_n}{\eta_n} \to 0$ gwarantuje separację skal czasowych. Dzięki temu proces $W_n$ zbiega asymptotycznie szybciej i z perspektywy wolniejszej iteracji $J_n$ może być traktowany jako quasi-stacjonarny.
\subsection{Lematy i twierdzenia}
\begin{lemma}[Kontrakcja operatora Bellmana dla QH]
\label{thm:bellman-contraction}
Definiujemy operator $T^{\phi_s}: \mathbb{R}^{|S|} \to \mathbb{R}^{|S|}$:
\begin{equation}
(T^{\phi_s} J)(s) = \sum_a \phi_s(a\mid s) \left[ r(s, a) + \sum_{s'} q(s'\mid s, a) \left[ -(1 - \alpha)\beta \sum_{a'} \phi_s(a'\mid s') r(s', a') + \beta J(s') \right] \right].
\end{equation}
Wtedy operator $T^{\phi_s}$ jest kontrakcją w normie $\max$ ze współczynnikiem dyskontowania $\beta < 1$. Istnieje jedyny punkt stały $J = T^{\phi_s} J$, $J$=$J^{\alpha, \beta}_{\phi_s}$.
\end{lemma}
\begin{proof}
Pokażemy, że $T^{\phi_s}$ jest kontrakcją w normie $\|\cdot\|_\infty$.
Niech $J_1, J_2 \in \mathbb{R}^{|S|}$ będą dowolne. Dla ustalonego stanu $s$ rozważmy różnicę
 $\bigl|(T^{\phi_s}J_1)(s) - (T^{\phi_s}J_2)(s)\bigr|$.
Korzystając z definicji operatora $T^{\phi_s}$ oraz faktu, że wyrazy niezależne od $J$ redukują się w różnicy, otrzymujemy
\[
\begin{aligned}
&(T^{\phi_s}J_1)(s) - (T^{\phi_s}J_2)(s) \\
&= \sum_a \phi_s(a\mid s) \sum_{s'} q(s'\mid s,a)\, \beta\bigl(J_1(s') - J_2(s')\bigr).
\end{aligned}
\]
Stąd
\[
\begin{aligned}
\bigl|(T^{\phi_s}J_1)(s) - (T^{\phi_s}J_2)(s)\bigr|
&\leq \beta \sum_a \phi_s(a\mid s) \sum_{s'} q(s'\mid s,a)\, \bigl|J_1(s') - J_2(s')\bigr| \\
&\leq \beta \sum_a \phi_s(a\mid s) \sum_{s'} q(s'\mid s,a)\, \|J_1 - J_2\|_\infty \\
&= \beta\,\|J_1 - J_2\|_\infty,
\end{aligned}
\]
ponieważ $\sum_a \phi_s(a\mid s)=1$ oraz $\sum_{s'} q(s'\mid s,a)=1$.
Maksymalizując po $s \in S$, dostajemy
\[
\|T^{\phi_s}J_1 - T^{\phi_s}J_2\|_\infty \leq \beta\,\|J_1 - J_2\|_\infty.
\]
Zatem dla $\beta \in [0,1)$ operator $T^{\phi_s}$ jest kontrakcją ze współczynnikiem $\beta$.
Z twierdzenia Banacha o punkcie stałym wynika istnienie i jednoznaczność punktu stałego $J^*$ takiego, że $T^{\phi_s}J^* = J^*$. Ponieważ równanie $J = T^{\phi_s}J$ jest równaniem Bellmana dla wartości QH przy polityce $\phi_s$, tym punktem stałym jest $J^{\alpha,\beta}_{\phi_s}$.
\end{proof}

\subsubsection{Zbieżność algorytmu}

\begin{theorem}[Zbieżność Algorytmu~\ref{alg:policy-eval}]
\label{thm:policy-eval-conv}
Jeżeli kroki uczenia $\eta_n$, $\theta_n$ spełniają warunki Robbins--Monro:
\[
\sum_n \eta_n = \infty, \quad \sum_n \eta_n^2 < \infty \quad \text{oraz} \quad \sum_n \theta_n = \infty, \quad \sum_n \theta_n^2 < \infty,
\]
 oraz dodatkowo $\frac{\theta_n}{\eta_n} \to 0$ (dwie skale czasowe), to ciągi $(W_n, J_n)$ są zbieżne prawie na pewno do $(J^{\alpha, \beta}_{\phi_s}, J^{\alpha, \beta}_{\mu, \phi_s})$.
\end{theorem}

Dowód Twierdzenia ~\ref{thm:policy-eval-conv} jest umieszczony w pracy \cite{eshwar2025}, jest on oparty o stochastyczną aproksymację \cite{borkar2008}
% \begin{proof}[Szkic dowodu]
% Dowód opiera się na teorii aproksymacji stochastycznej o dwóch skalach czasowych.

% Wprowadźmy operator oczekiwanej aktualizacji dla $W$ oraz $J$ wynikający z równania~\eqref{eq:final-policy-eval}. Przy założeniu pokrycia ($\nu(a\mid s)>0$ tam, gdzie $\mu$ lub $\phi_s$ wybierają akcje) aproksymatory oparte o wagi importance sampling mają skończoną wartość oczekiwaną i są nieobciążone.

% 	\textbf{Szybka skala ($W_n$):} Ponieważ $\theta_n/\eta_n\to 0$, iteracja $W_n$ widzi $J_n$ jako quasi-stałe i zachowuje się jak standardowa aproksymacja stochastyczna do punktu stałego odpowiadającego wartości $J^{\alpha,\beta}_{\phi_s}$.

% 	\textbf{Wolna skala ($J_n$):} W granicy $W_n \approx J^{\alpha,\beta}_{\phi_s}$, iteracja $J_n$ jest aproksymacją stochastyczną do równania liniowego z punktem stałym $J^{\alpha,\beta}_{\mu,\phi_s}$.

% Formalizacja przebiega standardowo przez metodę ODE dla aproksymacji stochastycznej z dwoma krokami uczenia i wykorzystuje fakt, że odpowiadające układy deterministyczne mają jednoznaczne globalnie asymptotycznie stabilne punkty stałe.
% \end{proof}

\section{Znajdowanie optymalnej polityki}

W tej sekcji zdefiniujemy czym jest optymalna polityka w kontekście dyskontowania quasi-hiperbolicznego oraz przedstawimy algorytm uczenia się optymalnej polityki bez znajomości środowiska modelu.
\subsection{Definicja optymalnej polityki}
Optymalna polityka w dyskontowaniu QH to para optymalna, która składa się z reguły decyzyjnej dla pierwszego kroku i polityki stacjonarnej $(\mu^\star, \phi_s^\star)$ maksymalizująca wartość funkcji wartości $J^{\alpha, \beta}_{\mu, \phi_s}(s)$ dla każdego stanu $s \in S$. Formalnie:
\begin{equation}
(\mu^\star, \phi_s^\star) = \arg\max_{\mu \in \Phi_m} \max_{\phi_s \in \Phi_s} J^{\alpha, \beta}_{\mu, \phi_s}(s), \quad \forall s \in S.
\end{equation}

Zatem aby uzyskać uzyskać optymalną politykę, wystarczy rozważać $(\mu, \phi_s)$ zgodnie z Twierdzeniem~\ref{thm:policy-reduction}.

$J^{\alpha, \beta}_{\mu^\star, \phi_s^\star}(s)$ = $\max\limits_{\mu,\phi_s} J^{\alpha, \beta}_{\mu, \phi_s}(s)$ = $\max\limits_{\mu,\phi_s} \sum\limits_a \mu(a\mid s) \left[ r(s, a) + \alpha\beta \sum\limits_{s'} q(s'\mid s, a) J^{\beta}_{\phi_s}(s') \right]$=\\= $\max\limits_{\mu} \sum\limits_a \mu(a\mid s) \left[ r(s, a) + \alpha\beta \sum\limits_{s'} q(s'\mid s, a) \max\limits_{\phi_s} J^{\beta}_{\phi_s}(s') \right]$. \\
Jako iż $\max\limits_{\phi_s} J^{\beta}_{\phi_s}(s)$ odpowiada maksymalnej wartości funkcji celu z dyskontowaniem wykładniczym, to zapisujemy optymalną politykę jako $\phi^\star_\beta$. Konsekwentnie, optymalna polityka dla dyskontowania QH jest dana przez 
\[
\phi_s^\star = \phi^\star_\beta,
\]
oraz
\begin{equation}
\mu^\star(s) = \arg\max_{\mu}\sum\limits_{a}\mu(a\mid s) \left[ r(s, a) + \alpha \beta \sum_{s'} q(s'\mid s, a) J^{\beta}_{\phi^\star_\beta}(s') \right].
\end{equation}
\subsection{Twierdzenie o istnieniu optymalnej polityki deterministycznej}

\begin{theorem}[Deterministyczna optymalna polityka]
\label{thm:deterministic-optimal}
Istnieje deterministyczna optymalna para, która składa się z reguły decyzyjnej dla pierwszego kroku i polityki stacjonarnej $(\mu^\star, \phi_s^\star)$ taka, że
\begin{equation}
V^{\alpha, \beta}_{\mu^\star, \phi_s^\star}(s) = \max_{\mu} \max_{\phi_s} V^{\alpha, \beta}_{\mu, \phi_s}(s), \quad \forall s \in S.
\end{equation}
Zauważmy, że para $(\mu^\star, \phi_s^\star)$ jest polityką deterministyczną markowską w MDP.
\end{theorem}
\begin{proof}
Najpierw musimy pokazać, że $\phi_s^\star$ może być wybrana jako deterministyczna. Dla dyskontowania wykładniczego z rozdziału 62.4 w \cite{puterman2014} optymalna polityka deterministyczna $\phi_\beta^\star$. Zatem skoro ustaliliśmy, że $\phi_s^\star = \phi_\beta^\star$, to $\phi_s^\star$ jest deterministyczna. \\
Następnym krokiem jest pokazanie, że $\mu^\star*$ również jest deterministyczna. Z definicji $\mu^\star$ mamy
\[\mu^\star(s) = \arg\max_{\mu}\sum\limits_{a}\mu(a\mid s) \left[ r(s, a) + \alpha \beta \sum_{s'} q(s'\mid s, a) J^{\beta}_{\phi^\star_s}(s') \right].\]
Polityka $\phi_s^\star$ jest deterministyczna, zatem funkcja wartości $J^\beta_{\phi_s^\star}(s)$ jest dobrze zdefiniowana. Wyrażenie wewnątrz maksymalizacji jest liniową funkcją od $\mu$, a skoro $\mu$ jest rozkładem prawdopodobieństwa, to maksymalizacja jest osiągana przez wybór akcji $a^\star$ maksymalizującej wyrażenie w nawiasie. Zatem $\mu^\star$ może być wybrana jako deterministyczna, wybierając akcję $a^\star$ z $q(a^\star \mid s)=1$ dla każdego stanu $s$. \\
W ten sposób pokazaliśmy, że zarówno $\phi_s^\star$, jak i $\mu^\star$ mogą być wybrane jako polityki deterministyczne, co kończy dowód.
\end{proof}

\subsection{Algorytm 2: QH Q-Learning}

Drugi algorytm pozwala na znalezienie optymalnej polityki bez znajomości środowiska modelu. Tzn. znajdujemy optymalną parę polityk $(\mu^\star, \phi_s^\star)$ maksymalizującą wartość funkcji wartości $J^{\alpha, \beta}_{\mu, \phi_s}(s)$ dla każdego stanu $s \in S$ bazując jedynie na próbkach doświadczeń z interakcji ze środowiskiem, nie znając przejść $q$ ani funkcji nagród $r$.\\
Do tego celu definiujemy funkcję Q dla danej pary polityk $(\mu, \phi_s)$, która spełnia właściwość:
\begin{equation}
Q^{\alpha, \beta}_{\mu, \phi_s}(s,a) = (1-\alpha) r(s,a) + \alpha Q^{\beta}_{\mu,\phi_s}(s,a),
\end{equation}
Długo terminowy komponent funkcji Q, $Q^{\beta}_{\mu,\phi_s}(s,a)$ nie zależy od $\mu$, zatem możemy zapisać:
\[
Q^{\beta}_{\mu,\phi_s}(s,a) = Q^{\beta}_{\phi_s}(s,a).
\]
Dla optymalnej pary polityk $(\mu^\star, \phi_s^\star)$, zapisujemy optymalną funkcję Q jako:
\begin{equation}
\begin{aligned}
Q^{\alpha,\beta}_{\mu^\star,\phi_s^\star}(s,a)
&= \max_{\mu,\phi_s}\Bigl[(1-\alpha)\,r(s,a) + \alpha\,Q^{\beta}_{\mu,\phi_s}(s,a)\Bigr] =\\
&=  \max\limits_{\phi_s}[(1-\alpha)r(s,a) + \alpha Q^\beta_{\phi_s}(s,a)] \\
&= (1-\alpha)\,r(s,a) + \alpha\,\max_{\phi_s} Q^{\beta}_{\phi_s}(s,a).
\end{aligned}
\end{equation}
 Z poprzedniej sekcji wiemy, że $\max\limits_{\phi_s} Q^{\beta}_{\phi_s}(s,a)$ odpowiada optymalnej polityce dla dyskontowania wykładniczego Q. Oznaczmy tę optymalną politykę jako $\phi^\star_\beta$. Wtedy optymalna polityka dla dyskontowania QH jest dana przez:
\[ \phi_s^\star = \phi^\star_\beta, \]
oraz
\[
\mu^\star(s) = \arg\max\limits_a Q^{\alpha, \beta}_{\mu^\star, \phi_s^\star}(s,a) = \arg\max\limits_a \left[ (1-\alpha) r(s,a) + \alpha Q^{\beta}_{\phi^\star_\beta}(s,a) \right].
\]
\subsubsection{Opis algorytmu}

\begin{algorithm}[H]
\caption{QH Q-Learning (Bezmodelowy)}
\label{alg:qh-qlearning}
\begin{algorithmic}[1]
\STATE \textbf{Wymagania:} Kroki uczenia $(\eta_n)$, $(\theta_n)$; współczynniki dyskontowania $\alpha$, $\beta$
\STATE \textbf{Inicjalizacja:} $W_0, Q_0 \in \mathbb{R}^{|S| \times |A|}$
\FOR{$n = 0, 1, 2, \ldots$}
\FOR{$(s, a) \in S \times A$}
\STATE Próbkuj $s' \sim q(\cdot\mid s, a)$ i zaobserwuj $r(s, a)$
\STATE $W_{n+1}(s, a) \leftarrow W_n(s, a) + \eta_n \left[ r(s, a) + \beta \max_{a'} W_n(s', a') - W_n(s, a) \right]$
\STATE $Q_{n+1}(s, a) \leftarrow Q_n(s, a) + \theta_n \left[ (1 - \alpha) r(s, a) + \alpha W_n(s, a) - Q_n(s, a) \right]$
\ENDFOR
\ENDFOR
\STATE \textbf{Zwraca:} $Q_n$ (QH) oraz $W_n$ (wykładnicze); 
\[ \mu^\star(s) = \arg\max\limits_a Q_n(s, a), \]
\[ \phi_s^\star(s) = \arg\max\limits_a W_n(s, a). \]
\end{algorithmic}
\end{algorithm}

Żeby estymować funkcję Q-wartości zarówno dla dyskontowania wykładniczego ($W$), jak i dla dyskontowania quasi-hiperbolicznego ($Q$), Algorytm~\ref{alg:qh-qlearning} utrzymuje dwie oddzielne tablice wartości, $W_n$ i $Q_n$. \\
Aktualizacje $W_n$ są wykonywane z krokiem uczenia $\eta_n$ i odpowiadają klasycznemu Q-learningowi dla dyskontowania wykładniczego z współczynnikiem $\beta$. Aktualizacje $Q_n$ są wykonywane z krokiem uczenia $\theta_n$ i wykorzystują bieżące wartości $W_n$ do obliczenia celu TD, zgodnie z definicją funkcji Q dla dyskontowania quasi-hiperbolicznego.\\
\subsubsection{Konwergencja algorytmu}

\begin{theorem}[Konwergencja QH Q-Learning]
\label{thm:qh-qlearning-conv}
Załóżmy, że kroki uczenia $(\eta_n)$ i $(\theta_n)$ spełniają warunki Robbins--Monro:
\[
\sum_n \eta_n = \infty, \quad \sum_n \eta_n^2 < \infty, \quad \sum_n \theta_n = \infty, \quad \sum_n \theta_n^2 < \infty,
\]
 oraz dodatkowo $\theta_n / \eta_n \to 0$ (dwie skale czasowe). 

Wtedy ciągi $(Q_n, W_n)$ generowane przez Algorytm QH Q-Learning zbiegają prawie na pewno do $(Q^{\alpha, \beta}_{\mu^\star, \phi_s^\star}, Q^{\beta}_{\phi_s^\star})$. 

Polityki ekstrahowane jako $\mu^\star(s) = \arg\max_a Q_n(s,a)$ oraz $\phi_s^\star(s) = \arg\max_a W_n(s,a)$ są optymalne.
\end{theorem}

\begin{proof}[Szkic dowodu]
Dowód wykorzystuje teorię aproksymacji stochastycznej o dwóch skalach czasowych:

  extbf{Szybka skala ($W_n$):}
Aktualizacja $W_n$ to klasyczny Q-learning dla dyskontowania wykładniczego:
\[
W_{n+1}(s, a) \leftarrow W_n(s, a) + \eta_n \left[ r(s, a) + \beta \max_{a'} W_n(s', a') - W_n(s, a) \right].
\]
Z klasycznej teorii Q-learningu wynika, że $W_n \to Q^{\beta}_{\phi_s^\star}$ prawie na pewno.

  extbf{Wolna skala ($Q_n$):}
Aktualizacja $Q_n$ ma postać:
\[
Q_{n+1}(s, a) \leftarrow Q_n(s, a) + \theta_n \left[ (1 - \alpha) r(s, a) + \alpha W_n(s, a) - Q_n(s, a) \right].
\]
Ponieważ $\theta_n / \eta_n \to 0$, proces $W_n$ działa w szybszej skali czasowej i można go traktować jako quasi-stacjonarny z punktu widzenia $Q_n$.

W granicy $W_n \approx Q^{\beta}_{\phi_s^\star}$, więc aktualizacja $Q_n$ zbieżna jest do układu liniowego ze stałym punktem $Q^{\alpha, \beta}_{\mu^\star, \phi_s^\star}$.

  extbf{Zachowanie greedy:}
Polityki greedy względem zbieżnych funkcji Q są optymalne:
\begin{align}
\mu^\star(s) &= \arg\max_a Q^{\alpha, \beta}_{\mu^\star, \phi_s^\star}(s,a), \\
\phi_s^\star(s) &= \arg\max_a Q^{\beta}_{\phi_s^\star}(s,a).
\end{align}

  Formalne uzasadnienie jest analogiczne do szkicu dowodu Twierdzenia~\ref{thm:policy-eval-conv}: również opiera się na metodzie ODE oraz analizie dwuskalowej. Szczegóły można znaleźć w pracy \cite{eshwar2024}.
\end{proof}

\subsubsection{Złożoność obliczeniowa}

Analizujemy złożoność czasową i pamięciową algorytmów:
\begin{itemize}
\item \textbf{Pamięć:} $O(|S| \times |A|)$ dla tablic $W$ i $Q$
\item \textbf{Czas na iterację:} $O(|S| \times |A|)$ dla pełnej aktualizacji
\item Algorytmy dwuskalowe wymagają większej liczby iteracji dla zbieżności, ale oferują gwarancje konwergencji
\end{itemize}

\subsection{Porównanie z metodami tradycyjnymi}

Algorytm QH Q-Learning różni się od klasycznego Q-learningu w następujących aspektach:

\begin{itemize}
\item \textbf{Dwie funkcje Q:} QH Q-Learning utrzymuje dwie funkcje Q: $W$ dla dyskontowania wykładniczego i $Q$ dla dyskontowania quasi-hiperbolicznego.
\item \textbf{Dwie skale czasowe:} Aktualizacje $W$ i $Q$ odbywają się z różnymi krokami uczenia, gdzie $\theta_n / \eta_n \to 0$.
\item \textbf{Polityka dwuskładnikowa:} Wynikiem jest para optymalna, która składa się z reguły decyzyjnej dla pierwszego kroku i polityki stacjonarnej $(\mu^\star, \phi_s^\star)$ zamiast pojedynczej polityki.
\item \textbf{Niespójność czasowa:} Optymalna polityka nie jest spójna w czasie, co wymaga specjalnej interpretacji i zastosowania.
\end{itemize}
